\documentclass[12pt,a4paper]{article}
\usepackage[utf8]{inputenc}
\usepackage{amsmath}
\usepackage{amssymb}
\usepackage{geometry}

\geometry{margin=1in}

\title{Derivation of MSE Gradient for Backpropagation}
\author{Neural Network Library}
\date{}

\begin{document}

\maketitle

\section{Introduction}

This document derives the gradient of the Mean Squared Error (MSE) loss function, which is used to initialize the backpropagation algorithm in the neural network library.

\section{Mean Squared Error Definition}

For a single training sample with output $y$ and target $t$, the squared error is:

\begin{equation}
E = (y - t)^2
\end{equation}

For $n$ training samples, the Mean Squared Error is:

\begin{equation}
MSE = \frac{1}{n} \sum_{i=1}^{n} (y_i - t_i)^2
\end{equation}

\section{Derivative with Respect to Output}

To perform backpropagation, we need the gradient of the MSE with respect to the network output. For a single sample:

\begin{align}
E &= (y - t)^2 \\
\frac{\partial E}{\partial y} &= \frac{\partial}{\partial y} (y - t)^2
\end{align}

Using the chain rule, let $u = y - t$:

\begin{align}
\frac{\partial}{\partial y} u^2 &= 2u \cdot \frac{\partial u}{\partial y} \\
&= 2(y - t) \cdot \frac{\partial}{\partial y}(y - t) \\
&= 2(y - t) \cdot 1 \\
&= 2(y - t)
\end{align}

\section{Scaling Factor for Mean Squared Error}

For the mean squared error averaged over $n$ samples:

\begin{align}
MSE &= \frac{1}{n} (y - t)^2 \\
\frac{\partial MSE}{\partial y} &= \frac{\partial}{\partial y} \left[ \frac{1}{n} (y - t)^2 \right] \\
&= \frac{1}{n} \cdot 2(y - t) \\
&= \frac{2}{n}(y - t)
\end{align}

\section{Connection to Code}

In \texttt{nn.h} at line 420 of \texttt{nn\_backprop\_traditional}:

\begin{verbatim}
MAT_AT(NN_OUTPUT(g), 0, j) = (MAT_AT(NN_OUTPUT(nn), 0, j) - MAT_AT(to, i, j))*2/n;
\end{verbatim}

This directly implements the derived formula:

\begin{equation}
\frac{\partial MSE}{\partial y_j^{(i)}} = \frac{2}{n} \cdot (y_j^{(i)} - t_j^{(i)})
\end{equation}

Where:
\begin{itemize}
    \item $y_j^{(i)}$ = \texttt{MAT\_AT(NN\_OUTPUT(nn), 0, j)} — network output
    \item $t_j^{(i)}$ = \texttt{MAT\_AT(to, i, j)} — target value
    \item The result is stored in the gradient structure \texttt{NN\_OUTPUT(g)}
\end{itemize}

\section{Summary}

\begin{center}
\begin{tabular}{|c|c|}
\hline
\textbf{Mathematical Notation} & \textbf{Code Implementation} \\
\hline
$\frac{\partial MSE}{\partial y}$ & \texttt{NN\_OUTPUT(g)} \\
\hline
$y$ (output) & \texttt{NN\_OUTPUT(nn)} \\
\hline
$t$ (target) & \texttt{to} \\
\hline
$n$ (samples) & \texttt{n} \\
\hline
\end{tabular}
\end{center}

\end{document}
